\section{Evidence-Gated Expansion of the Frontier}
\label{sec:ood}\label{sec:frontier}

\begin{designbox}
Model capability growth as evidence-gated, not automatic. A candidate joins the
retained set only when it verifies against accepted evidence; the union notation
captures that intent, not a promise. In practice the set is revisable and not
empirically monotone, and a run that adds no capability can still add durable value
by recording what not to repeat.
\end{designbox}

The two preceding sections give a fixed-model state $H$ that missions update
(Section~\ref{sec:evolution}) and a retained process that outlives each artifact
(Section~\ref{sec:process}). Their intended effect on the outermost spine stage ---
the expanded out-of-distribution frontier --- is an evidence-gated capability set.
Let $C_t$ be the capabilities the system can bring to a new, out-of-distribution task
at time $t$, and $\epsilon$ the acceptance threshold of the evidence gate.

\begin{equation}
C_{t+1}
= C_t \cup
\{c:\operatorname{Verify}(c,E_t)\ge\epsilon\}.
\end{equation}

Read precisely, the equation models one thing: the \emph{retained verified additions}.
A candidate capability $c$ --- a new skill, a synthesized wiki technique, a validated
procedure --- enters the set only when $\operatorname{Verify}(c,E_t)$ clears the gate
that the Reviewer owns (Section~\ref{sec:rolesplanes}). The union is deliberately
conservative about \emph{intent}: the frontier is meant to grow by verified increments
rather than by unchecked accumulation.

\subsection{What the notation does not claim}

The set notation is an idealization, and the report is explicit about its limits, in
line with the master spine's own caption that capability is not guaranteed to grow
every run. It \textbf{does not guarantee} that any given mission succeeds, that
practical capability is empirically \textbf{monotone}, or that an admitted capability
is retained forever. Real $C_t$ is revisable: skills and wiki pages can be revised,
archived, or retired through the ordinary lifecycle (reinforce, distill, revise,
retire), so an entry can leave the effective set, and a capability verified once can
regress on a harder instance. The clean union is a statement of design intent, not an
empirical monotonicity claim.

Expansion also has a mode that adds no new $c$ at all. A mission that reaches a
\textbf{negative result} or a \textsc{no-go} decision still writes durable evidence
about what does not work; recorded in the process (Section~\ref{sec:process}), it can
reduce the repeated search cost of later missions, which is real progress even when
$C_{t+1}=C_t$. Finally, two expansion directions are distinct and should not be
conflated: \emph{breadth} across domains --- taking on a new arena or vertical --- and
\emph{depth} within a domain --- solving harder instances of a problem already in
range. The public results and portfolio that follow are read against this qualified
frontier, not against a promise of universal or monotone improvement
(Sections~\ref{sec:results}, \ref{sec:portfolio}, \ref{sec:limitations}).
